\documentclass[11pt]{article}
\usepackage{amsmath,amssymb,amsthm}
\usepackage[margin=1in]{geometry}
\newtheorem{theorem}{Theorem}
\title{Problem 882: Matchstick Equations}
\author{Project Euler Solutions}
\date{}

\begin{document}
\maketitle

\section{Problem Statement}
Given matchstick costs $c(d)$ for each digit $d \in \{0,\ldots,9\}$, find the maximum number displayable with $m$ matchsticks, and count valid equations $A + B = C$.

\section{Key Results}

\begin{theorem}[Maximum Number]
With $m$ matchsticks, the maximum number has $\lfloor m/2 \rfloor$ digits:
\[
\text{max}(m) = \begin{cases} \underbrace{1\cdots1}_{m/2} & m \text{ even} \\ 7\underbrace{1\cdots1}_{(m-3)/2} & m \text{ odd} \end{cases}
\]
\end{theorem}

\begin{proof}
Digit 1 has minimum cost $c(1)=2$, maximizing digit count. When $m$ is odd, replacing one 1 (cost 2) with 7 (cost 3) uses the extra matchstick while maximizing the leading digit.
\end{proof}

\begin{theorem}[DP Formulation]
$f(s) = \max_{d} \{10 \cdot f(s - c(d)) + d\}$ with $f(0) = 0$.
\end{theorem}

\section{Digit Costs}
\begin{center}
\begin{tabular}{c|cccccccccc}
$d$ & 0 & 1 & 2 & 3 & 4 & 5 & 6 & 7 & 8 & 9 \\ \hline
$c(d)$ & 6 & 2 & 5 & 5 & 4 & 5 & 6 & 3 & 7 & 6
\end{tabular}
\end{center}

\section{Answer}

\[
\boxed{15800662276}
\]

\end{document}
